\documentclass[10pt]{article}
\usepackage[margin=20mm]{geometry}
\usepackage{amsmath,amssymb,fontspec,unicode-math,xurl,hyperref}
\IfFontExistsTF{DejaVu Serif}{\setmainfont{DejaVu Serif}}{\setmainfont{Latin Modern Roman}}
\setmathfont{latinmodern-math.otf}
\setlength{\parindent}{0pt}\setlength{\parskip}{6pt}\sloppy
\title{Selected evidence and evaluation checks for EMMA}
\author{AI-assisted technical review}
\date{}
\begin{document}
\maketitle
Abstract: A claim-focused review checked the benchmark composition, selected score aggregations, the stated test-time selection comparison, evaluation-source logic, and a visual electrostatics example. The independent calculations agree with the inspected principal summaries. No substantive defect is established in that scope. This is not a certification of the 2,788 answer keys, historical model outputs, or a universal necessity for visual reasoning. Adjudication is pending.\par
\section*{Identity and reviewed material}
This report reviews Yunzhuo Hao, Jiawei Gu, Huichen Will Wang, Linjie Li, Zhengyuan Yang, Lijuan Wang and Yu Cheng, Can MLLMs Reason in Multimodality? EMMA: An Enhanced MultiModal ReAsoning Benchmark, ICML 2025, PMLR 267, OpenReview `v26vwjxOEz`. The frozen proceedings PDF has 47 pages. The main text and the applicable benchmark/protocol appendices B–C were read, with selected breakdown tables. The extensive case-study appendix was not exhaustively solved or visually inspected. The release at https://github.com/EMMA-Bench/EMMA, commit `47259952dd8e95bc1e3855a4164ccd3634f8eb1e`, was inspected statically for query construction, answer extraction, equality testing and aggregation. The acquired snapshot excludes large data/result JSON files; it does not permit an item-level reproduction of the historical scores.\par
\section*{Concrete checks and their outcomes}
\subsection*{Composition and score weighting}
Table 2, page 6, uses subject counts \((892,156,1176,564)\), summing to 2,788. Overall full-benchmark accuracy is an item-weighted quantity, \[A=\frac{892A_{\rm math}+156A_{\rm phys}+1176A_{\rm chem}
                 +564A_{\rm code}}{2788},\] not the unweighted average of four subject percentages. Using the printed subject rates gives \(35.0788\%\) for direct Claude 3.5 Sonnet, \(32.4247\%\) for direct GPT-4o, and \(42.5039\%\) for Gemini Thinking-0121. These agree with the corresponding printed \(35.08\%\), \(32.42\%\) and \(42.50\%\) at the precision justified by rounded inputs. EMMA-mini has 100 questions per subject, so its overall score is instead the equally weighted subject average. The human values \((75,64.5,86,85.5)\) average to \(77.75\%\). The gap to \(48\%\) is 29.75 percentage points, as stated.\par
The coding construction in B.2, pages 16–17, yields \(47(4+4+2+2)=564\) questions, agreeing with the task counts \((188,188,94,94)\). It also explains why coding category membership is nonexclusive. Summing category scores as if they partitioned the dataset would be inappropriate, but the inspected accumulator adds each question only once to the overall and subject totals while allowing multiple coding category memberships.\par
\subsection*{Scaling results and error taxonomy}
Table 3, page 7, distinguishes actual selection methods from Pass@\(N\) oracle coverage. The best reported selection score, \(50.75\%\), remains 27 percentage points below the human average; Pass@16 of \(81.5\%\) is an upper bound requiring knowledge of which response is correct. These statements are consistent: high oracle coverage does not establish that the tested selector can identify a correct response. The results support the bounded statement that the evaluated methods, models and budgets did not close the gap. They do not rule out larger budgets or different selectors. The paper reports different decoding temperatures for one-sample and multi-sample experiments in Table 6, page 20; this prevents reading the comparison as a controlled estimate of sampling count alone.\par
For the 106 incorrect o1 answers in the selected math/coding analysis, the displayed percentages in Figure 5, page 8, correspond to counts \((32,56,10,8)\): perceptual, visual-reasoning, textual-reasoning and knowledge errors. The resulting percentages are \((30.1887,52.8302,9.4340,7.5472)\%\), agreeing with the figure. This validates the arithmetic, not the annotators’ attribution of each individual error. The two-expert-per-subject average is a reported human baseline, not a population-level estimate established by this review.\par
\subsection*{An independently solved source image}
Figure 1, page 1, depicts charges \(+Q,-2Q,+3Q\) at the vertices of an equilateral triangle and asks for the force on the top \(+Q\). Taking side length one and positions \[r_+=(0,\sqrt3/2),\quad r_-=(-1/2,0),\quad r_3=(1/2,0),\] with the common factor \(kQ^2\) suppressed, the two forces are \((-1,-\sqrt3)\) and \((-3/2,3\sqrt3/2)\). Hence \[F=(-5/2,\sqrt3/2),\] pointing \(19.1066\) degrees above the leftward horizontal, consistent with the figure’s answer b and its correction of the model’s repulsion direction. The relevant source image was actually viewed. Figure 6 and its current-loop diagram on page 8 were also visually inspected for the error-analysis context, but no separate quantitative magnetic-field validation is claimed.\par
\section*{Supporting checks}
The filtering protocol in Section 3.2, page 4, removes an item if one of three models solves its text-plus-caption version at least five times out of ten. This is an operational, model- and caption-dependent selection rule. It is not a proof that no sufficiently detailed textual description or future text reasoner can solve a retained problem. For perspective, under an illustrative independent four-option guessing model, the chance of at least five successes in ten trials is \[\sum_{k=5}^{10}\binom{10}{k}(1/4)^k(3/4)^{10-k}=0.0781269.\] This is a property of the toy binomial model, not an estimate of the actual filter’s error rate; correlations, nonuniform option counts and open-ended questions preclude such an extrapolation. Manual review and the observed performance gap provide additional evidence, but do not turn the filtering rule into a theorem about modality necessity.\par
The inspected evaluator accepts direct letters, boxed answers and selected final-answer phrases, then tests strings, number words and mathematical equivalence, including a numerical two-decimal comparison. It also offers an optional model-based evaluation mode. These are real released-code behaviors; the exact mode and output files used for each historical result have not been established. No API calls or author scripts were executed. Corpus-level answer-key validity, chemical structures, parsing failure rates, contamination, expert labeling agreement and the independence of generated candidates remain outside the inspected evidence. The main text’s retained-chemistry count of 42 differs from the appendix’s 20; the appendix’s \(892+80+20=992\) retained total reconciles the headline figure. This is a provenance clarification requiring the actual curation manifest, not evidence here of an incorrect scored answer or material result change.\par
\section*{Conclusion}
No substantive issue is established in the checked aggregation, selected scaling comparisons, source logic and electrostatics example. The conclusion is deliberately narrower than full-benchmark validation.\par
Dimension | State\par
Inspection assessment | No Issue Found within the declared scope\par
Repair verification | Not applicable; no correction proposed\par
Scientific adjudication | Pending; initial review only\par
Author adoption | Unknown; no contact made\par
Remaining work | Item-level data/output audit, additional examples and independent acceptance\par
\section*{Evidence and review limits}
The source is the frozen PMLR proceedings PDF, identified by SHA-256\par
> `c904d6272c272eb72acc5641ea3e56b2ed672a1c4641feb052fc628ad5b2705b`\par
Physical PDF page numbers are one based and coincide with the paper’s printed page numbers in this version. The check program was written for this review; no author program was executed. Published training runs, data collection and reported performance were not reproduced. It is not independent acceptance. No human review or author adoption is recorded. Scientific adjudication remains pending.\par
9 Yunzhuo Hao; Jiawei Gu; Huichen Wang; Linjie Li; Zhengyuan Yang; Lijuan Wang; Yu Cheng. Can MLLMs Reason in Multimodality? EMMA: An Enhanced MultiModal ReAsoning Benchmark. Proceedings of the 42nd International Conference on Machine Learning, PMLR 267 (2025). Published proceedings PDF, frozen for this review on 25 October 2026. https://proceedings.mlr.press/v267/hao25a.html.\par
\end{document}
